\chapter*{Preface}
\addcontentsline{toc}{chapter}{Preface}
% Front matter, on the operator's directives (2026-07-25 13:48 + 13:57 +
% 14:04-05): the preface IS the photoelectric effect, static friction, and
% stress-strain curves (the three phenomena opening, + water hammer + the
% GPS temporal-friction weave per the operator's unifying chain "fourth
% satellite IS static friction IS photoelectric IS water hammer"); then earn
% the number 4 as the solution to the binary result of trilateration; then
% the 4 models. Laws held: each phenomenon cited (Amontons 1699/Coulomb
% 1785; Lenard 1902/Einstein 1905; Hooke 1678/Young 1807/Cauchy 1822;
% Joukowsky 1898); the SHARED SHAPE is front-matter framing (the series'
% own image, "one shape of response," never one mechanism); trilateration
% classical; NS owed in one clause; model-fence echoed; no Lean. Gate: Kodo
% (citation) + Beastmaster (assembled architecture re-pass).

Push a block across a floor, gently, and then less gently. For a while
nothing moves at all: the force climbs and the block sits, absorbing the
push without an inch of response. Then, at a threshold, it lets go all at
once. The laws of this stubbornness are old --- Amontons wrote them down in
1699, and Coulomb in 1785 separated the sticking from the sliding --- and
every reader has felt them through a heavy piece of furniture: the response
of the block is not proportional to the push. It is nothing, and then it is
everything.

Shine light on a metal and ask for electrons. Brighten a red lamp as much as
you please --- pour energy in --- and nothing comes out, ever; switch to a
dim violet one and the electrons leave at once. The observation is Lenard's,
in 1902; the explanation is Einstein's, in 1905, and it carried the century:
the metal does not respond to the amount of light accumulated but to whether
each arriving quantum clears a threshold. Below the boundary, nothing; above
it, release --- and no patience with the red lamp substitutes for crossing.

Pull a bar of steel and plot the pull against the stretch. At first the plot
is a clean straight line --- the response proportional to the load, the
regime Hooke announced in 1678 and Young made quantitative in 1807, the
regime where the material meets force with obedience. Then the line reaches
the yield point and stops being a line: the proportionality ends at a
boundary, and past it the bar answers by a different law entirely. The
formal language of that plot --- stress against strain --- is Cauchy's, from
1822, and its shape is the same two-act story: a linear zone, and an edge
where the linear zone runs out.

Close a valve on a long pipe of moving water. Close it slowly and nothing of
note happens. Close it fast and the pipe bangs like a struck bell --- the
water hammer, whose surge Joukowsky measured and wrote into an equation in
1898. The moving column does not negotiate with the closed valve; its answer
arrives whole, as a wave, at the moment the boundary is imposed.

Four phenomena, four mechanisms, and no claim here that they share one
physics --- each has its own, cited above, and this book will not blur them.
What they share is a \emph{shape of response}, and the reader should hold
that shape, because this series is organized by it: a zone where response is
proportional and accumulation buys progress; a boundary where proportion
ends; and, at the boundary, a response that is withheld and then arrives
whole. Nothing, then everything. That shape --- offered here as this
preface's own framing image, nothing more --- is the one the reader's pocket
already demonstrates in time. A receiver listening to three satellites does
not close in gradually on its position; it sticks, holding two candidate
answers it cannot choose between, and it waits --- time accumulating against
the ambiguity like force against a stuck block, a temporal friction ---
until the fourth reading arrives, or a valid calibration stands in for it,
and the fix releases, whole. Why exactly four, and why the sticking is
structural rather than sloppy, is the arithmetic this preface now performs.

So before the construction begins, a number has to be earned, because this
book is about to be organized by it.

How many soundings does it take to fix a point? The question is older than
any instrument in this series; it is the surveyor's question, and its
arithmetic is classical. One measured distance confines the unknown to a
sphere around the station. A second distance cuts that sphere with another,
and the intersection of two spheres is a circle. A third cuts again, and a
circle meets a sphere in exactly two points. There the classical count
stalls, and it stalls in an interesting way: three honest measurements,
performed perfectly, leave not a haze of candidates but precisely two --- a
point and its mirror twin, utterly indistinguishable to the three stations
that measured them. This is the binary result of trilateration, and it is
structural, not sloppy: no repetition of the same three soundings, however
careful, will ever split the pair, because the ambiguity lives in the
geometry, not in the noise.

The solution is the fourth sounding. A fourth station, placed off the plane
of the first three, is near one twin and far from the other; its single
reading breaks the tie, and the point, at last, has been named. The
arithmetic of the earning is worth saying aloud: three dimensions demand
three ranges to come down to finitely many candidates, and the finitely many
is exactly two --- a binary digit of ambiguity --- so one more measurement,
the smallest possible act, decides it. Four is three plus the tiebreaker.
The count is not a convention and not a preference; it is the smallest
number of soundings that turns ranging into naming, in a world of three
dimensions. The reader carries the living example in a pocket: satellite
navigation fixes a receiver from four satellites --- the classical geometry
just told, with the honest engineering footnote that the fourth reading is
also made to pay for the receiver's imperfect clock, a second job that
changes nothing about the first.

The number is now earned, and here is what it organizes. Across this series,
four models are being developed --- four descriptions ranged on the same
unknown. Volume 1 opened its preface with a fence this book keeps: the
electron, in these pages, is the name of a mathematical model, and the
argument touches descriptions, never the particle in the world. It is
exactly because the electron is a model that it can be sounded the way a
surveyor sounds a point: each description is a station, and no single
station names the thing. There is the geometric description, in which the
model's content appears as curvature; the quantum description, in which it
appears as states and their duals; the electron model itself, the device's
own subject, the one description with a measured number to its name; and the
fluid description --- promised in the founding text of the apparatus before
any of the others were built, and, at this writing, still the promise the
series has yet to pay. Four stations. Three would leave the position
ambiguous --- a mirror pair of readings with no way to choose --- and the
whole series is, in this exact sense, a trilateration: it does not expect
any single description to name its subject, and it does not stall at three.

This volume builds none of the four. It builds the ground they stand on ---
the common language of neighborhoods, charts, tensors, connections, and
curvature in which such descriptions are written at all --- and it builds
that language entirely from cited mathematics, name and date attached, with
nothing original anywhere in the construction. A surveyor does not invent
trigonometry; a surveyor learns it, cites it, and points the instrument.
That is this book's whole posture. The four soundings belong to the volumes
around it; the mathematics belongs to the names in the back matter; and the
reader, as in Volume 1, is owed this notice on the first page: what follows
asks them to believe nothing new. It only asks them to watch the classical
instruments assemble.
